\section{Complementary statements for the actual return record}
\label{app:retained-return-statements}
The statements below collect the return-record conclusions and the
intermediate inversion formulas used in the article. Their hypotheses
and proofs are unchanged. Theorem~\ref{thm:intro-stationary-microscopic}
now evaluates the complete stationary physical probability in Theorem X;
the four-coordinate return-density problem remains a distinct endpoint.
The letters A--X retain the cross-reference convention used in the
geometric and probabilistic sections.

\paragraph{The physical record.}
Let $T_R$ be the collision map, $Y_R^*$ the section constructed below,
and $F_R^*$ its actual first-return map. Write
$J_{n,R}=(K_{n,R},N_{n,R},T_{n,R})$ for the four-coordinate record of
$n$ returns, with the two coordinates of $K_{n,R}$ displayed separately.
Its law is taken under the normalized collision measure $\nu_R^*$.
The following synopsis records the unconditional probabilistic
conclusions. The geometry, normalization, and limiting matrices are
specified intrinsically in the cited theorems.

\begin{maintheorem}[Gaussian laws of the physical record]
\label{thm:intro-gaussian}
For $R\in I=[0.45,0.47]$ set
\[
 c_* = \frac{91}{10000\pi},\qquad
 \bar\tau_R=\frac{\sqrt3/2-\pi R^2}{2R},\qquad
 \bar G_R=(0,0,c_*^{-1},\bar\tau_R/c_*).
\]
There is a continuous positive semidefinite matrix $D_R$ such that,
for every $K<\infty$ and $n\ge2$,
\[
 \sup_{R\in I,\ |t|\le K}
 \left|\int e^{it\cdot(J_{n,R}-n\bar G_R)/\sqrt n}\dd\nu_R^*
              -e^{-t^{\mathsf T}D_Rt/2}\right|
 \le C_Kn^{-1/26}\sqrt{\log(2+n)}.
\]
The polygonal centered return processes converge, uniformly in $R$,
to Brownian motion with covariance $D_R$. The stationary physical
processes of lattice displacement and centered collision count also
satisfy a parameter-uniform functional Gaussian limit.

More precisely, $D_R=c_*^{-1}\Gamma_R$, where $\Gamma_R$ is the
absolutely convergent collision Green--Kubo matrix of
$h_R=f_R-\bar G_R\one_{Y_R^*}$. Its kernel consists exactly of directions
$v$ for which $v\cdot(G_R-\bar G_R)$ is an $L^2$ coboundary for
$F_R^*$. No nondegeneracy assumption is used in these conclusions.
\end{maintheorem}
\begin{proof}
Theorems~\ref{thm:collision-covariance},
\ref{thm:actual-record-gaussian}, \ref{thm:gaussian-kernel-coboundary},
and \ref{thm:actual-return-functional}, together with
Corollary~\ref{cor:physical-functional-gaussian}, prove the assertions.
\end{proof}

\begin{maintheorem}[An integrated central band]
\label{thm:intro-major-arc}
For the same actual return record and the matrix $D_R$ in Theorem A,
\[
 \sup_{R\in I}\int_{|v|\le2n^{1/200}}
 \left|\int e^{iv\cdot(J_{n,R}-n\bar G_R)/\sqrt n}\dd\nu_R^*
                       -e^{-v^{\mathsf T}D_Rv/2}\right|\dd v
       \le Cn^{-3/280}\sqrt{\log(2+n)}.
\]
The variable $v$ is rescaled: the corresponding physical frequencies
have size at most $2n^{-1/2+1/200}$. The same estimate holds for complex
initial insertions whose zero extensions to the collision cylinder
have uniformly bounded supremum and variation norms, with the Gaussian
multiplied by the integral of the insertion.
\end{maintheorem}
\begin{proof}
Apply Theorem~\ref{thm:growing-major-arc}; the original probability is
the admissible insertion $a_R=s_R$.
\end{proof}

\paragraph{Complex insertion convention.}
In Theorem B the insertion is a density with respect to the collision
probability $\nu$. Its amplitude $\alpha_R=\int a_R\dd\nu$ may be
complex; the associated pushforward is a finite complex measure, not
in general a probability. The original section probability corresponds
to $a_R=s_R=c_*^{-1}\one_{Y_R^*}$.

\begin{maintheorem}[A mark at any actual return]
\label{thm:intro-marked-return}
Let $a$ be a complex function on the collision cylinder, zero outside
$Y_R^*$, and put $M_a=\|a\|_\infty$, $V_a=\|a\|_{\mathrm{BV}}$,
and $\alpha_a=\int a\dd\nu$. For $0\le k\le n$ set
\[
 C^{[k],a}_{n,R}(v)=\int_{Y_R^*}a((F_R^*)^k x)
       e^{iv\cdot(J_{n,R}(x)-n\bar G_R)/\sqrt n}\dd\nu(x).
\]
There is a constant $C$, independent of $R,n,k,a$, such that, for $n\ge2$,
\begin{align*}
 \int_{|v|\le2n^{1/200}}
 \left|C^{[k],a}_{n,R}(v)-\alpha_a e^{-v^{\mathsf T}D_Rv/2}\right|\dd v
 \le C\left(M_an^{-3/280}\sqrt{\log(2+n)}+V_an^{-9/175}\right).
\end{align*}
Thus the same central band accommodates initial, terminal, and one
intermediate return-state insertion, including a mark whose index
varies with $n$. The displayed transform is integrated against the
unnormalized restricted collision measure $\nu|_{Y_R^*}$; division by
$c_*$ gives the corresponding transform under the section probability
$\nu_R^*$. No positivity of $a$ or nonsingularity of $D_R$ is assumed.
\end{maintheorem}
\begin{proof}
This is Theorem~\ref{thm:marked-return-band}.
\end{proof}


\begin{maintheorem}[Actual second moments and a marked quadratic law]
\label{thm:intro-actual-moments}
Put $U_{n,R}=J_{n,R}-n\bar G_R$. Uniformly for $R\in I$,
\[
 \int |U_{n,R}|^4\dd\nu_R^*\le Cn^2,\qquad
 \left\|\frac1n\int U_{n,R}\otimes U_{n,R}\dd\nu_R^*-D_R\right\|
       \le Cn^{-1/6}.
\]
The count component satisfies $e_3^{\mathsf T}D_Re_3\ge d_N>0$
uniformly in $R$, where $e_3=(0,0,1,0)$.
For the complex insertion and normalization of Theorem C,
\begin{align*}
 &\left\|\int_{Y_R^*}a((F_R^*)^kx)
       \frac{U_{n,R}(x)\otimes U_{n,R}(x)}n\dd\nu(x)-\alpha_aD_R\right\|\\
 &\hspace{20mm}\le C\left(M_an^{-1/6}+(M_a+V_a)n^{-1}\right),
       \qquad 0\le k\le n.
\end{align*}
The matrix $D_R$ is also the uniform limit of the induced Green--Kubo
formula with Cesaro weights. No induced mixing or absolute convergence
of its correlation series is assumed.
\end{maintheorem}
\begin{proof}
Apply Theorems~\ref{thm:marked-L2-stopping} and
\ref{thm:marked-quadratic-moments}, and
Corollaries~\ref{cor:induced-cesaro-covariance} and
\ref{cor:positive-count-variance}.
\end{proof}

\begin{maintheorem}[Uniform nondegeneracy of the joint record]
\label{thm:intro-joint-nondegeneracy}
The covariance of the same actual four-coordinate return record satisfies
\[
 d_0 I_4\le D_R\le D_0 I_4\qquad(R\in I)
\]
for constants $0<d_0\le D_0<\infty$. Thus every nonzero real joint
direction has positive Gaussian variance. For all sufficiently large
$n$, the normalized actual covariance in Theorem D is at least
$(d_0/2)I_4$, uniformly in $R$. The three-dimensional physical-time
fluctuation covariance is uniformly positive definite as well.
\end{maintheorem}
\begin{proof}
Apply Theorem~\ref{thm:joint-uniform-nondegeneracy} and
Corollary~\ref{cor:actual-uniform-ellipticity}.
\end{proof}
In every subsequent use of the matrix from Theorems A--D, the stronger
uniform positive definiteness of Theorem E is available. The positive
semidefinite formulation of the earlier statements records their weaker
logical input, not a remaining covariance hypothesis.


\begin{maintheorem}[Physical phase arithmetic and quantitative defects]
\label{thm:intro-quantitative-phases}
Let $\omega=(u,s,t)\in\T^2\times\T\times\R$ and
$\phi_{R,\omega}=u\cdot\kappa_R+s+t\tau_R$ on the collision cylinder.
A measurable circle-valued solution of
$q\circ T_R=e^{i\phi_{R,\omega}}q$ exists exactly when $\omega=0$;
in that case it is constant almost everywhere.

For each fixed $R$ and compact nonzero frequency set $\mathcal K$,
there is $c_{R,\mathcal K}>0$ such that, for $H\ge2$,
\[
 \inf_{\omega\in\mathcal K}
 \|q\circ T_R-e^{i\phi_{R,\omega}}q\|_1
                  \ge\frac{c_{R,\mathcal K}}{1+\log H}
 \quad (|q|=1,\ \|q\|_{\mathrm{BV}}\le H).
\]
For complex $f$ with $\|f\|_2=1$ and
$\|f\|_\infty+\|f\|_{\mathrm{BV}}\le H$, the corresponding bound is
\[
 \inf_{\omega\in\mathcal K}
 \|f\circ T_R-e^{i\phi_{R,\omega}}f\|_2
                  \ge\frac{c_{R,\mathcal K}}{(1+\log H)^2}.
\]
No lower bound on $|f|$ is assumed. For fixed $H$ and $\mathcal K$,
the second infimum is also positive uniformly over $R\in I$.
The latter assertion does not specify a uniform logarithmic rate in $H$.
\end{maintheorem}
\begin{proof}
Theorem~\ref{thm:complete-physical-phase} proves the exact assertion.
Theorems~\ref{thm:compact-log-phase-defect},
\ref{thm:compact-vector-defect} and
\ref{thm:uniform-compact-bv-separation} prove the defect bounds.
\end{proof}

\begin{maintheorem}[Near-origin defects for the actual return record]
\label{thm:intro-actual-return-defects}
Set $g_R=G_R-\bar G_R$ and
\[
 \Lambda=1+\log(H/r),\qquad K=\Lambda\log(2+\Lambda),
                 \qquad H\ge2,\quad 0<r<1/2.
\]
There are $a,c,r_*>0$, independent of $R$, such that $r\le r_*$
and $rK\le a$ imply the following conclusions for $|z|=r$.
Every unit-modulus section function $q$ whose zero extension has
variation at most $H$ satisfies
\[
 \|q\circ F_R^*-e^{iz\cdot g_R}q\|_{L^1(\nu_R^*)}\ge cr^2.
\]
Every complex section function $f$ with
$\|f\|_{L^2(\nu_R^*)}=1$ and
$\|f\|_\infty+\|f^0\|_{\mathrm{BV}}\le H$ satisfies
\[
 \|f\circ F_R^*-e^{iz\cdot g_R}f\|_{L^2(\nu_R^*)}
                                    \ge cr^2/K.
\]
The function may vanish. With $H\le H_0n^\zeta$, these bounds apply
uniformly on the physical annulus
$2n^{-99/200}\le|z|\le n^{-2/5}$ for all sufficiently large $n$.
The corresponding rescaled radii are $2n^{1/200}$ and $n^{1/10}$.
\end{maintheorem}
\begin{proof}
Apply Theorems~\ref{thm:actual-return-circle-defect} and
\ref{thm:actual-return-vector-defect}, and
Corollary~\ref{cor:actual-annulus-defect}.
\end{proof}
\paragraph{The spectral parameter in Theorem G.}
The displayed defect in Theorem G has unit spectral parameter for the
centered return multiplier. A constant per return is not the same as a
constant per collision. The next theorem treats both the physical
frequency $z$ and the return spectral phase $\xi$, with
$\lambda=e^{i\xi}$, and strengthens the regularity constraint.

\begin{maintheorem}[A joint peripheral neighborhood]
\label{thm:intro-peripheral-neighborhood}
For $\xi\in[-\pi,\pi]$, set
\[
 \rho=(|z|^2+\xi^2)^{1/2},\quad
 \Lambda=1+\log(H/\rho),\quad K=\Lambda\log(2+\Lambda).
\]
There are $a,b,\rho_0>0$, independent of $R$, such that
$0<\rho\le\rho_0$ and $\rho^2K\le a$ imply
\[
 \|q\circ F_R^*-e^{i(\xi+z\cdot g_R)}q\|_{L^1(\nu_R^*)}
       \ge b\rho^2
\]
for every unit section phase with $\|q^0\|_{\mathrm{BV}}\le H$,
and
\[
 \|f\circ F_R^*-e^{i(\xi+z\cdot g_R)}f\|_{L^2(\nu_R^*)}
       \ge b\rho^2/K
\]
for every normalized complex section function with
$\|f\|_\infty+\|f^0\|_{\mathrm{BV}}\le H$. Zeros are allowed.
These estimates include $z=0$, $\xi\ne0$. They are local in the joint
parameters, not estimates over the full peripheral circle.

For the physical annulus $2n^{-99/200}\le|z|\le n^{-2/5}$,
$|\xi|\le n^{-\gamma}$, and any
$0<\beta<\min\{4/5,2\gamma\}$, they apply to regularity budgets
$H\le H_0\exp(n^\beta)$ for all sufficiently large $n$.
The rescaled annulus is $2n^{1/200}\le|\sqrt n\,z|\le n^{1/10}$.
\end{maintheorem}
\begin{proof}
Apply Theorem~\ref{thm:joint-return-spectral-defect} and
Corollary~\ref{cor:peripheral-annulus-budgets}.
\end{proof}

Theorem~\ref{thm:compressed-peripheral-resolvent} converts these physical
defects into a lower singular-value bound for finite-rank orthogonal
compressions of the actual return operator. The reconstruction has
supremum and variation budget $C/h$ at mesh $h$; no positive lower
modulus is assumed. Proposition~\ref{prop:compressed-finite-time-comparison}
separately bounds the difference between the compressed powers and the
unmodified characteristic function. A mesh-uniform spectral limit and
the complete complementary-frequency integral are not inferred from
that finite-time comparison.


\begin{maintheorem}[Finite-count extraction of the raw density]
\label{thm:intro-finite-count-extraction}
Fix $R\in I$, $n\ge1$ and a collision-count cutoff $L$. For every
bounded finite-record subanalytic weight in
Definition~\ref{def:finite-record-weight}, the restricted actual law
\[
 \mu_{n,R}^{w,L}=(J_{n,R})_*
                    (w\one_{\{N_{n,R}\le L\}}\nu_R^*)
\]
has an exact decomposition $\mu_{n,R}^{w,L}=E^L+Q^L$.
The extracted density is a finite sum of cutoff one-sided terms
$x^\alpha(\log x)^k$, with $\alpha\in\mathbb Q$, $-1<\alpha\le1$,
and $k$ a nonnegative integer. It includes all singular boundary
contributions to the finite-count law. Each lattice component of the
residual density $r^L$ belongs to $W^{2,1}(\R)$. In particular, if
$A_2=\sum_\ell\|\partial_t^2 r^L(\ell,\cdot)\|_1$, then
\[
 \widehat Q^L\in L^1(\T^3\times\R),\qquad
 \int_{\T^3\times\{|b|>B\}}|\widehat Q^L|\dd\omega
                   \le 2(2\pi)^3 A_2/B .
\]
Here $A_2$ is finite at each fixed $(n,L,R,w)$; no uniform long-time
bound for this quantity is part of the assertion.

For the original unweighted law, choose $L_n=\lceil\lambda n\rceil$
with $\lambda>c_*^{-1}+1$. Its density agrees exactly with the restricted
density on every fixed central window for large $n$.
The finite extraction gives the raw inversion identity
\eqref{eq:count-localized-exact-inversion} on that window. The
high-count measure enters only through a low-frequency convolution,
whose normalized central contribution is $O(n^{-P})$ for every $P>0$.
The physical cutoff is $2n^{-99/200}$, or $2n^{1/200}$ after rescaling.
\end{maintheorem}
\begin{proof}
Apply Theorem~\ref{thm:finite-count-raw-extraction},
Lemma~\ref{lem:central-count-support}, and
Theorem~\ref{thm:finite-extraction-central-inversion}.
\end{proof}

Theorem I supplies the complete finite-count residual, rather than only
an initial-coordinate variation bound. Corollary~\ref{cor:finite-extraction-error-budget}
displays the remaining uniform residual-growth, complementary-integral
and local-edge estimates for the original raw density.


\begin{maintheorem}[Exponential reconstruction and multiple observations]
\label{thm:intro-window-reconstruction}
The finite-record variation budget may be taken as $Ce^{CL}$.
Consequently, with $\rho=(|z|^2+\xi^2)^{1/2}$ and
$\Lambda=1+\log(H/\rho)$, the joint return spectral estimates of
Theorem H strengthen to
\[
 \|\mathcal E_{R,z,\xi}q\|_1\ge b\rho^2,\qquad
 \|\mathcal E_{R,z,\xi}f\|_2\ge b\rho^2/\Lambda
\]
whenever $0<\rho\le\rho_0$ and $\rho^2\Lambda\le a$.
Here $q$ is a unit section phase of zero-extension variation at most
$H$, while $f$ is a normalized complex section function with supremum
plus zero-extension variation at most $H$. The constants are uniform
in $R$, and no positive lower modulus of $f$ is required.

There are constants $a_*,\theta_*,\beta_*>0$ for the original family
with the following property. Let $m_n\le a_*\log n$ and choose any
$0\le\ell_1<\cdots<\ell_q\le m_n$. Let $E_j\subset Y_R^*$ be
observation sets whose zero-extended indicators have total variation
at most a fixed power of $\log(2+n)$. If the unchanged event
\[
 A_{n,k,R}=\bigcap_{j=1}^{q}
       \{x:(F_R^*)^{k+\ell_j}x\in E_j\},\qquad
                       0\le k\le n-m_n,
\]
has probability at least $p_0 n^{-\beta_*}$, then, for
$U_{n,R}=J_{n,R}-n\bar G_R$,
\[
 \sup_{R,k}\int_{|v|\le2n^{\theta_*}}
 \left|\mathbb E[e^{iv\cdot U_{n,R}/\sqrt n}\mid A_{n,k,R}]
                          -e^{-v^{\mathsf T}D_Rv/2}\right|\dd v
                           \longrightarrow0.
\]
Its conditional normalized mean tends to zero and its conditional
covariance tends to $D_R$, uniformly. All expectations use $\nu_R^*$.
The physical cutoff is $2n^{-1/2+\theta_*}$. Explicit constants and
more general probability, variation and frequency tradeoffs are given
in \eqref{eq:window-exponent-margins} and
\eqref{eq:nonempty-window-budget}.
\end{maintheorem}
\begin{proof}
Theorems~\ref{thm:exponential-finite-record} and
\ref{thm:single-log-return-defect} prove the first part.
Corollary~\ref{cor:logarithmic-window-conditioning} and
Theorem~\ref{thm:window-weighted-moments}, with
\eqref{eq:nonempty-window-budget}, prove the remaining assertions.
\end{proof}

The single-logarithm estimate is local in the joint physical and return
spectral parameters; its exact lifting identity remains valid at every
return phase. The growing observation window is a new class, not an
assumed BV bound on a long pullback. Its smaller positive central band
is additional to the unchanged $n^{1/200}$ band of Theorem C.
Neither statement replaces the complementary-integral or local-edge
estimates for the raw density. The original multiple-time event is
never exchanged for an event with a collision cutoff.

\begin{maintheorem}[Direct annular averages for the uncompressed record]
\label{thm:intro-direct-annular-averages}
Let $\mathcal U_{R,z}f=e^{-iz\cdot g_R}f\circ F_R^*$ on
$L^2(\nu_R^*)$. For all sufficiently large $n$, all $T\ge n$,
$N\ge0$, $a\in L^2(\nu_R^*)$, and physical frequencies
\[
 2n^{-99/200}\le|z|\le n^{-2/5},
 \qquad 2n^{1/200}\le|\sqrt n\,z|\le n^{1/10},
\]
one has, uniformly in $R$,
\begin{align*}
 \frac1T\sum_{j=0}^{T-1}
 \left|\int a e^{-iz\cdot U_{N+j,R}}\dd\nu_R^*\right|^2
                          &\le C n^{-4/495}\|a\|_2^2,\\
 \sup_{\xi\in\R}
 \left\|\frac1T\sum_{j=0}^{T-1}
         e^{-ij\xi}\mathcal U_{R,z}^{N+j}\one\right\|_2^2
                          &\le C n^{-4/495}.
\end{align*}
For an unchanged event $A$ of probability $p>0$, the average of the
squared conditional characteristic functions is at most
$C p^{-1}n^{-4/495}$. No BV condition on $\one_A$ is needed.

If $\varsigma_{R,ru}$ is the spectral probability of $\one$ for
$\mathcal U_{R,ru}$, where $|u|=1$, its principal angle divided by
$r^2$ converges in law, uniformly in $R,u$, to the Cauchy law of scale
$u^{\mathsf T}D_Ru/2$.
\end{maintheorem}
\begin{proof}
Theorems~\ref{thm:direct-moving-averages} and
\ref{thm:weighted-annular-averages}, with
Corollary~\ref{cor:direct-annular-averages}, prove the averaged bounds.
Theorem~\ref{thm:spectral-Cauchy-limit} proves the spectral conclusion.
\end{proof}

\paragraph{The prescribed-count target.}
Theorem K concerns averages in the return count. The raw local theorem
instead requires, at the prescribed $n$, an estimate of the form
$n^2\int|1-\chi_n|\,|H_{n,R}|\dd\omega=o(1)$ for the appropriate
edge-subtracted residual, together with its local extracted-edge term.
The averaged bound does not imply this estimate: its existing annular
normalization gives only $n^{2/5-2/495}$ after Cauchy--Schwarz and the
raw factor $n^2$.

The common lag scale is $m(|z|)=\lfloor|z|^{-200/99}\rfloor<n/4$.
It requires no spatial mesh, and no BV reconstruction of a time-$n$
vector. The estimates are time averages or cyclic spectral statements,
not decay at each specified return count. The physical record still has
the Gaussian law in Theorem A; the Cauchy law here concerns a different,
explicitly identified spectral probability.
The Cauchy characteristic convergence is for fixed rescaled spectral
test frequency. It supplies no uniform estimate at $t=n|z|^2$
when that quantity grows across the physical annulus.

\begin{maintheorem}[Fixed-return cancellation and a larger central band]
\label{thm:intro-fixed-return-annulus}
For the insertion and measure conventions of Theorem C, let
$\Phi^{[k],a}_{n,R}(z)=\int_{Y_R^*}a((F_R^*)^kx)e^{iz\cdot J_{n,R}(x)}\dd\nu(x)$.
Uniformly for $R\in I$ and $0\le k\le n$,
\begin{align*}
 &n^2\int_{2n^{-99/200}\le|z|\le2n^{-12/25}}
                    |\Phi^{[k],a}_{n,R}(z)|\dd z\\
 &\qquad\le C\left(M_an^{-1/200}\sqrt{\log(2+n)}+V_an^{-13/100}\right).
\end{align*}
This holds at every prescribed sufficiently large return count, not
only on average. Moreover
\begin{align*}
 &\int_{|v|\le2n^{1/50}}
 |C^{[k],a}_{n,R}(v)-\alpha_a e^{-v^{\mathsf T}D_Rv/2}|\dd v\\
 &\qquad\le C\left(M_an^{-1/200}\sqrt{\log(2+n)}+V_an^{-9/175}\right).
\end{align*}
The new physical support radius is $2n^{-12/25}$; the old,
sharper estimate on $2n^{-99/200}$ remains unchanged. The unweighted
law is $a=s_R$. Theorem~\ref{thm:expanded-count-localized-inversion}
places the enlarged cutoff in the exact raw-density decomposition.
\end{maintheorem}
\begin{proof}
Apply Theorem~\ref{thm:fixed-return-inner-annulus} and
Corollary~\ref{cor:fixed-count-expanded-band}.
\end{proof}

\begin{maintheorem}[Damped unsmoothing and a wider fixed-return band]
\label{thm:intro-damped-unsmoothing}
For the insertion, measure and frequency conventions of Theorems C and L,
\begin{align*}
 &n^2\int_{2n^{-99/200}\le|z|\le2n^{-19/42}}
          |\Phi^{[k],a}_{n,R}(z)|\dd z
        \le C\left(M_an^{-5/861}+V_an^{-59/21}\right),\\
 &\int_{|v|\le2n^{1/21}}
   \left|C^{[k],a}_{n,R}(v)-\alpha_a e^{-v^{\mathsf T}D_Rv/2}\right|\dd v
        \le C\left(M_an^{-5/861}+V_an^{-9/175}\right).
\end{align*}
Both estimates hold at every sufficiently large prescribed return count,
uniformly in $R$, $0\le k\le n$ and the stated bounded-BV insertions.
The physical support radius is $2n^{-19/42}$ and its rescaled radius is
$2n^{1/21}$. More generally, a fixed-order construction reaches every
rescaled exponent $1/200<\varepsilon<1/20$ with a positive decay rate.
The prior central and inner-annulus theorems remain valid on their bands.
\end{maintheorem}
\begin{proof}
Theorems~\ref{thm:damped-cubic-unsmoothing} and
\ref{thm:wider-fixed-count-annulus},
Corollary~\ref{cor:wider-marked-central-band}, and
Proposition~\ref{prop:damped-unsmoothing-band-range} prove the assertions.
\end{proof}

The new weighted raw identity in
Theorem~\ref{thm:sharp-weighted-raw-inversion} uses this larger cutoff
and recomputes its local edge correction. Its finite-band residual,
long-time derivative budget and exact-event comparison are not replaced
by the new marked characteristic estimate.


\begin{maintheorem}[Multiscale stopping and a rate-preserving band]
\label{thm:intro-multiscale-stopping}
For the actual return record and insertion convention of Theorems C and D,
let $\mathcal G_d(D)$ denote the degree-$d$ moment tensor of $N(0,D)$.
For every fixed integer $d\ge1$, put $\rho_d=1/2$ for odd $d$ and
$\rho_d=1$ for even $d$. Then
\begin{align*}
 &\left\|\int_{Y_R^*}a((F_R^*)^kx)
       \left(\frac{U_{n,R}(x)}{\sqrt n}\right)^{\otimes d}\dd\nu(x)
                      -\alpha_a\mathcal G_d(D_R)\right\|\\
 &\hspace{15mm}\le C_d\left(M_an^{-1/4}+(M_a+V_a)n^{-\rho_d}\right).
\end{align*}
In particular, the actual normalized covariance and its induced Cesaro
formula converge to $D_R$ at rate $O(n^{-1/4})$. Moreover,
\begin{align*}
 &n^2\int_{2n^{-99/200}\le|z|\le2n^{-633/1400}}
                    |\Phi^{[k],a}_{n,R}(z)|\dd z
      \le C\left(M_an^{-3/280}+V_an^{-983/350}\right),\\
 &\int_{|v|\le2n^{67/1400}}
       |C^{[k],a}_{n,R}(v)-\alpha_a e^{-v^{\mathsf T}D_Rv/2}|\dd v\\
 &\hspace{15mm}\le C\left(M_an^{-3/280}\sqrt{\log(2+n)}
                                   +V_an^{-9/175}\right).
\end{align*}
The estimates are uniform in $R$, the prescribed sufficiently large
return count $n$, and one arbitrary mark $0\le k\le n$. The two support
radii are $2n^{-633/1400}$ physically and $2n^{67/1400}$ after rescaling.
No return-count averaging or normalization by annular volume is used.
The same-event consequence retains the original probability and variation
budget of Theorem C on this larger band.
\end{maintheorem}
\begin{proof}
Apply Theorems~\ref{thm:multiscale-marked-stopping},
\ref{thm:all-marked-gaussian-moments}, and
\ref{thm:rate-preserving-annulus}, together with
Corollary~\ref{cor:rate-preserving-same-event}.
\end{proof}

\begin{maintheorem}[Separate coordinate widths at the prescribed count]
\label{thm:intro-anisotropic-band}
Let
\[
 \mathcal B_n^{\rm roof}
   =[-2n^{1/100},2n^{1/100}]^3
                         \times[-2n^{9/100},2n^{9/100}].
\]
For the same insertion and measure convention as Theorem C,
\begin{align*}
 &n^2\int_{\substack{\sqrt n z\in\mathcal B_n^{\rm roof}\\
                            |z|\ge2n^{-99/200}}}
                  |\Phi^{[k],a}_{n,R}(z)|\dd z
       \le C\left(M_an^{-1/40}+V_an^{-497/25}\right).
\end{align*}
The estimate is uniform in $R$, one prescribed mark $0\le k\le n$,
and all sufficiently large return counts $n$. On the union
$\mathcal U_n=\mathcal B_n^{\rm roof}\cup\{|v|\le2n^{67/1400}\}$,
\begin{align*}
 &\int_{\mathcal U_n}
 |C^{[k],a}_{n,R}(v)-\alpha_a e^{-v^{\mathsf T}D_Rv/2}|\dd v\\
 &\hspace{12mm}\le C\left(M_an^{-3/280}\sqrt{\log(2+n)}
                                      +V_an^{-9/175}\right).
\end{align*}
The new physical widths are $2n^{-49/100}$ in the three discrete
coordinates and $2n^{-41/100}$ in flight time. They extend the proved
region without asserting cancellation on the whole isotropic annulus.
More generally separate rescaled exponents $\theta_i\ge1/200$ are
admissible when $\max_i\theta_i<1/10$ and
$\sum_i\theta_i+\max_i\theta_i<1/4$, with the rates specified in
Theorem~\ref{thm:anisotropic-fixed-count-band}.
\end{maintheorem}
\begin{proof}
Apply Theorem~\ref{thm:high-order-damped-unsmoothing},
Theorem~\ref{thm:anisotropic-fixed-count-band}, and
Corollary~\ref{cor:roof-box-fixed-count}.
\end{proof}

\begin{maintheorem}[A nonrectangular region at the original central rate]
\label{thm:intro-adaptive-cross}
Let $\mathcal H_n$ be the dyadic region defined in
\eqref{eq:adaptive-cross-scales}--\eqref{eq:adaptive-cross-support}.
Its base scale is $2n^{1/200}$, its maximal coordinate scale is
$n^{9/100}$, and its product-times-maximum budget is
$n^{67/280}/[\log(2+n)]^3$. Put
\[
 \mathcal U_n^\sharp=\mathcal H_n\cup\mathcal B_n^{\rm roof}
                         \cup\{|v|\le2n^{67/1400}\}.
\]
For the actual prescribed count, mark and measure convention of
Theorem C,
\begin{align*}
 \int_{\mathcal U_n^\sharp}
 |C^{[k],a}_{n,R}(v)-\alpha_a e^{-v^{\mathsf T}D_Rv/2}|\dd v
 \le C\left(M_an^{-3/280}\sqrt{\log(2+n)}+V_an^{-9/175}\right).
\end{align*}
The physical region is $n^{-1/2}\mathcal U_n^\sharp$. It retains
both the old isotropic ball and the roof box, and reaches physical
scale $n^{-41/100}$ on every coordinate axis. It is not an isotropic
ball of radius $n^{-2/5}$. Constants are uniform in $R,n,k,a$;
the residual order $P=26$ and spectral degree $Q=29$ are fixed.
\end{maintheorem}
\begin{proof}
Use Theorem~\ref{thm:adaptive-cross-fixed-count} and
Corollary~\ref{cor:adaptive-cross-directions-event}.
\end{proof}

\begin{maintheorem}[Coherent transport of the raw correction]
\label{thm:intro-coherent-raw-transport}
For an additionally finite-record admissible marked weight, fix one
exact finite-count extraction $\mu^{a,L}=E^{a,L}+Q^{a,L}$.
For a cutoff $\chi$ with kernel $K_\chi$, put
\[
 \mathcal R_\chi^{a,L}=e^{a,L}-K_\chi*E^{a,L}
            +\mathcal F^{-1}[(1-\chi)\widehat Q^{a,L}].
\]
For any two of the new dyadic, preceding box and preceding isotropic
cutoffs, and $L_n=\lceil\lambda n\rceil$ with
$\lambda>c_*^{-1}+1$, on each fixed central window,
\begin{align*}
 n^2\mathop{\rm ess\,sup}_{\mathcal W_{n,R,A}}
 |\mathcal R_\chi^{a,L_n}-\mathcal R_\psi^{a,L_n}|
 \le C\left(M_an^{-3/280}+V_an^{-983/350}\right)
        +C_{A,\lambda,P_0}M_an^{-P_0}
\end{align*}
for every fixed $P_0>0$. The difference is a bounded continuous
function even when an individual correction has infinite essential
supremum. This compares the entire signed correction, not either of
its two absolute bounds separately.
\end{maintheorem}
\begin{proof}
Proposition~\ref{prop:exact-raw-cutoff-transport} and
Theorem~\ref{thm:coherent-raw-cutoff-rate} prove the assertion.
\end{proof}

\begin{maintheorem}[Unsmoothed windows and the common signed correction]
\label{thm:intro-window-remainder}
Let $U_{n,R}=J_{n,R}-n\bar G_R$, $\rho=67/1400$, and
$B_n(\xi)=\prod_{i=1}^4[\xi_i-n^{-1/25},\xi_i+n^{-1/25}]$.
Uniformly on $|\xi|\le A$ and $R\in I$,
\[
 \nu_R^*\{U_{n,R}/\sqrt n\in B_n(\xi)\}
 =16n^{-4/25}g_{D_R}(\xi)
       \left(1+O_A(n^{-11/1400}\sqrt{\log(2+n)})\right).
\]
The event uses the unmodified record and actual integer intervals.
More generally, a rectangular frequency box contained in the retained
region gives the window estimates in
Theorem~\ref{thm:unsmoothed-window-denominator}, including positive
marked weights with their exact mass.

For a finite-record admissible complex mark, the same physical windows
$\mathcal B=n\bar G_R+\sqrt nB_n(\xi)$ satisfy
\[
 \frac{n^2}{\mathfrak m(\mathcal B)}
 \left|\int_{\mathcal B}\mathcal R_\chi^{a,L_n}\dd\mathfrak m\right|
 \le C_A\left(M_an^{-11/1400}\sqrt{\log(2+n)}+V_an^{-9/175}\right)
\]
for each of the three retained cutoffs and $L_n=\lceil\lambda n\rceil$,
$\lambda>c_*^{-1}+1$. Here $\mathfrak m$ is counting measure times
Lebesgue measure. The absolute value is outside the integral; the
assertion is not an essential-supremum bound.
\end{maintheorem}
\begin{proof}
Apply Corollary~\ref{cor:concrete-four-window} and
Theorem~\ref{thm:averaged-common-raw-remainder}.
\end{proof}

\begin{maintheorem}[Relative windows at deterministic physical time]
\label{thm:intro-relative-physical-window}
Start the orbit at a section collision with law $\nu_R^*$. Let
$V_R^Y(t)=(W_R^Y(t),C_R^Y(t)-t/\bar\tau_R)$ and
$n_R(t)=\lfloor c_*t/\bar\tau_R\rfloor$.
For $B_t=\prod_{i=1}^3[\xi_i-t^{-1/25},\xi_i+t^{-1/25}]$, define
\[
 A_t^{\rm phys}=\{V_R^Y(t)/\sqrt t\in B_t\},\qquad
 A_t^{\rm ret}=\{B_RU_{n_R(t),R}/\sqrt t\in B_t\},
\]
where $B_R(x)=(x_1,x_2,x_3-x_4/\bar\tau_R)$.
Uniformly for $|\xi|\le A$ and $R\in I$, both probabilities equal
\[
 8t^{-3/25}g_{\mathcal V_R}(\xi)
       \left(1+O_A(t^{-11/1400}\sqrt{\log(2+t)})\right),
\]
and
\[
 \frac{\nu_R^*(A_t^{\rm phys}\mathbin\triangle A_t^{\rm ret})}
      {\nu_R^*(A_t^{\rm ret})}
       \le C_A t^{-11/1400}\sqrt{\log(2+t)}.
\]
The conditional initial laws obey the same total variation bound.
These are exact physical window events of half-width $t^{23/50}$,
not singleton lattice events or events under a different initial law.
\end{maintheorem}
\begin{proof}
This is Theorem~\ref{thm:relative-physical-window}.
\end{proof}

\begin{maintheorem}[Stationary physical observations and conditional laws]
\label{thm:intro-stationary-physical-windows}
Represent equilibrium flow by the actual return suspension
\[
 \dd\mathbb P_R(y,s)=\frac{c_*}{\bar\tau_R}\dd\nu_R^*(y)\dd s,
 \qquad 0\le s<\tau_R^*(y).
\]
Let $A_t^{\rm phys}$ be the event that the stationary physical
increment $(W_R(t),C_R(t)-t/\bar\tau_R)/\sqrt t$ lies in the
three-dimensional box of center $\xi$ and half-width $t^{-2/25}$.
For $n_R(t)=\lfloor c_*t/\bar\tau_R\rfloor$, let $A_t^{\rm ret}$
be the event that $B_RU_{n_R(t),R}(y)/\sqrt t$ lies in the same box,
where $B_R(x_1,x_2,x_3,x_4)=(x_1,x_2,x_3-x_4/\bar\tau_R)$.
Then, uniformly in $R$ and bounded $\xi$,
\begin{align*}
 \mathbb P_R(A_t^{\rm phys})
 &=8t^{-6/25}g_{\mathcal V_R}(\xi)
       \left(1+O(t^{-23/2800}\sqrt{\log(2+t)})\right),\\
 \frac{\mathbb P_R(A_t^{\rm phys}\triangle A_t^{\rm ret})}
             {\mathbb P_R(A_t^{\rm ret})}
 &\le C t^{-23/2800}\sqrt{\log(2+t)}.
\end{align*}
The return event has the same denominator asymptotic. The same
comparison holds with the actual last-completed-return event, and
for the conditional laws of any common bounded path statistic.
The three physical coordinate half-widths are $t^{21/50}$.
The return origin $y$ has its equilibrium roof-length bias; it is
not drawn from the unbiased section probability.
\end{maintheorem}
\begin{proof}
Theorems~\ref{thm:wide-collision-band} and
\ref{thm:projected-collision-windows} provide the collision-clock
window estimate. Lemma~\ref{lem:stationary-length-bias} and
Proposition~\ref{prop:stationary-observation-coupling} give the
normalization and coupling. Theorem~\ref{thm:stationary-physical-windows}
proves the displayed assertions and their conditional-law consequence.
\end{proof}

\begin{maintheorem}[Physical paths conditioned on stationary windows]
\label{thm:intro-stationary-window-bridge}
Let $\mathsf X_{t,R}$ be a continuous interpolation of
\[
 s\longmapsto t^{-1/2}
   \bigl(W_R(ts),C_R(ts)-ts/\bar\tau_R\bigr),\qquad 0\le s\le1,
\]
whose supremum distance from the physical path is $O(t^{-1/2})$.
With the stationary initial law and the exact endpoint windows of
Theorem T, uniformly in $R$ and bounded $\xi$,
\[
 \operatorname{Law}(\mathsf X_{t,R}\mid A_t^{\rm phys})
       \ \Longrightarrow\
 \operatorname{Law}\bigl(s\xi+
     \mathcal V_R^{1/2}(W(s)-sW(1))\bigr)
       \quad\hbox{on }C([0,1],\R^3).
\]
Here $W$ is standard three-dimensional Brownian motion. Uniformity
is in bounded-Lipschitz distance. The same limit holds under the
prescribed-return, last-completed-return and collision reference events
on the same stationary space. For fixed interior times and bounded
boxes of nonempty interior, imposing those additional path observations
multiplies the endpoint denominator by their explicitly positive
Gaussian bridge probability; the resulting conditional path limit is
the bridge restricted to those boxes. The physical endpoint half-width
remains $t^{21/50}$.
\end{maintheorem}
\begin{proof}
Theorem~\ref{thm:multiblock-collision-band} and
Proposition~\ref{prop:window-conditional-characteristic} give joint
endpoint-window comparisons uniform in the block boundaries.
Theorem~\ref{thm:collision-window-bridge} proves conditional tightness.
Lemma~\ref{lem:whole-physical-clock-coupling},
Theorem~\ref{thm:stationary-physical-bridge} and
Corollary~\ref{cor:bridge-cylinder-events} give the physical conclusions.
\end{proof}

\paragraph{Microscopic reconstruction and conditional-measure reductions.}
Theorems V and W below are unconditional approximation theorems at the
original lattice resolution. They are separated from the Gaussian
asymptotic conclusions above: the remaining finite integral is not
identified with a Gaussian main term.

\begin{maintheorem}[Quantitative extraction and microscopic inversion]
\label{thm:intro-geometric-raw-reduction}
Use the normalized section probability and the exact finite-count
measure $\mu_{n,R}^{w,L}$. For a fixed product of bounded-BV return-state
marks, let $M,V$ be the budgets in \eqref{eq:geometric-weight-budgets}.
There is an exact decomposition $\mu_{n,R}^{w,L}=E_\varepsilon+Q_\varepsilon$
whose residual density $q_\varepsilon$ satisfies, uniformly in $R,n\le L$
and the mark locations,
\begin{align*}
 \|E_\varepsilon\|_{\TV}
    &\le C\{\varepsilon V+M(L+1)\varepsilon^{1/16}\},\\
 \sum_{\ell\in\Z^3}\|\partial_t^2q_\varepsilon(\ell,\cdot)\|_1
    &\le CM\exp\{C(L+1)\log(C/\varepsilon)\}.
\end{align*}
For $w=1$ both measures are positive. For polynomial $M,V$, any fixed
$P>0$ and any fixed upper bound on $|I|$, a roof cutoff with
$\log B_n=O(n\log n)$ gives
\[
 \sup_{R,\ell,I}n^2\left|
 \mu_{n,R}^w(\{\ell\}\times I)
       -\mathcal I_{n,R,B_n}^{w}(\ell,I)\right|\le C_Pn^{-P}.
\]
Here $\mathcal I$ is the finite Fourier integral of the full, unmodified
actual law in \eqref{eq:actual-microscopic-finite-integral}. No lattice
coordinate is averaged, and $I$ need not grow. This is a reduction of
the microscopic probability, not its Gaussian asymptotic or a pointwise
bound for the removed density.
\end{maintheorem}
\begin{proof}
Theorem~\ref{thm:geometric-extraction-budget} proves the extraction.
Proposition~\ref{prop:geometric-far-roof},
Theorem~\ref{thm:microscopic-finite-band-reduction} and
Corollary~\ref{cor:linear-cutoff-microscopic-inversion} give the inversion.
\end{proof}

\begin{maintheorem}[Positive microscopic reconstruction on the original space]
\label{thm:intro-positive-microscopic-conditioning}
Use either the normalized section law or the stationary return
suspension law \eqref{eq:equilibrium-source-for-inversion}. In the latter
case the record starts at the actual preceding section origin and all
trajectory statistics retain the original elapsed age.
For every $P>0$ there is a roof bandwidth $B_n$ with
$\log B_n=O(n\log n)$ and a nonnegative mixed density $p_{n,R,B_n}$
of mass one such that
\[
 n^2\|p_{n,R}-p_{n,R,B_n}\|_{L^1(\Z^3\times\R)}\le C_Pn^{-P}
\]
uniformly in $R$. The density is the finite integral
\eqref{eq:positive-full-law-density}; no lattice coordinate is smoothed.
For every microscopic event $A=\{J_{n,R}\in\{\ell\}\times I\}$, there
is a likelihood $0\le\lambda_{A,B_n}\le1$ on the same probability
space with
\[
 n^2\sup_{|W|\le1}
 \left|\mathbb E[W\one_A]-\mathbb E[W\lambda_{A,B_n}]\right|
        \le C_Pn^{-P}.
\]
Here $I$ is any bounded interval, including a shrinking interval, and
$W$ is any bounded measurable trajectory statistic. The second
expectation has the finite Fourier formula
\eqref{eq:all-selector-Fourier-formula}. More generally, a polynomial
number of intervals per lattice fibre is permitted. Normalization gives
a conditional total-variation bound after division by the same actual
denominator, as stated in Corollary~\ref{cor:positive-microscopic-posterior}.
\end{maintheorem}
\begin{proof}
Proposition~\ref{prop:stationary-geometric-extraction} supplies the
positive residual for both initial laws. Apply
Theorem~\ref{thm:uniform-positive-microscopic-inversion} and
Corollary~\ref{cor:positive-microscopic-posterior}.
\end{proof}

\begin{maintheorem}[Exact inversion and Gaussian-plus-complement reduction]
\label{thm:intro-exact-physical-endpoint}
Fix the same convex polygonal fundamental cell at both physical
observation times. For a prescribed physical collision count $m\ge1$,
let
\[
 p_{m,R}(k,t)=\mathbb P_R\{W_R(t)=k,\ C_R(t)=m\}.
\]
The initial and terminal within-flight cell corrections enter the exact
three-frequency inverse \eqref{eq:stationary-absolute-inversion}.
Uniformly in $m,R$, its time profiles satisfy
\[
 \sum_k\|D_t^2p_{m,R}(k,\cdot)\|_{\TV}\le C,
 \qquad
 \sum_k\|p_{m,R}(k,\cdot)-p^{\rm sharp}_{m,R,B}(k,\cdot)\|_\infty
                                                      \le C/B.
\]
The derivative is a signed measure. With $\Sigma_R$ and
$\zeta_{m,R}$ defined in
Section~\ref{sec:stationary-gaussian-complement}, the central physical
radius $2m^{-99/200}$, or $2m^{1/200}$ after rescaling, gives
\begin{align*}
 m^{3/2}p_{m,R}(k,t)
 &=\bar\tau_Rg_{\Sigma_R}(\zeta_{m,R}(k,t))
       +m^{3/2}\mathfrak M_{m,R,B}(k,t)\\
 &\quad+O\left(m^{-11/700}\sqrt{\log(2+m)}+m^{3/2}/B\right),
\end{align*}
where $\mathfrak M_{m,R,B}$ is the explicitly specified signed finite
middle-frequency inverse. Thus $B_m=m^{P+3/2}$ gives a pointwise
far-frequency error $O(m^{-P})$ on the normalized physical scale,
without a discarded-density term.

For the same original physical event there is a positive likelihood
$0\le\lambda_{m,k,B}\le1$ whose unnormalized source total-variation
error is at most $C/B$, uniformly against all bounded measurable path
selectors. A relative posterior estimate requires its actual selected
denominator. The microscopic Gaussian denominator requires the signed
middle term to vanish; it is not asserted by the tail bound alone.
\end{maintheorem}
\begin{proof}
Theorems~\ref{thm:stationary-overlap-regularity},
\ref{thm:stationary-polynomial-band},
\ref{thm:physical-time-selector-reconstruction} and
\ref{thm:physical-singleton-reduction} prove the statements.
\end{proof}

